
On sait (EGA II 8.1.8) que @@M0@@ est un sous-schéma fermé de @@M1@@ défini par un idéal de @@M2@@, canoniquement isomorphe à @@M3@@; l'immersion @@M4@@ étant définie par un idéal inversible, est une immersion fermée régulière de codimension @@M5@@. On note @@M6@@ l'idéal de @@M7@@ qui définit @@M8@@, et @@M9@@ l'idéal de @@M10@@ qui définit @@M11@@; on a donc @@M12@@. Enfin, on note @@M13@@, et @@M14@@ les morphismes structuraux. On a donc le diagramme commutatif:
@@M15@@

Par définition, @@M16@@, où @@M17@@. Par suite, on a:
@@M18@@
Or, @@M19@@ étant régulière, @@M20@@ (1.3). On note @@M21@@ le faisceau conormal @@M22@@; on a donc:
@@M23@@

Sur @@M24@@, on peut écrire la suite exacte:
@@M25@@

On remarque que @@M26@@; c'est donc le faisceau conormal de l'immersion @@M27@@, et on le notera @@M28@@.

\textbf{Lemme 3.2.} \uline{Avec les hypothèses et les notations de} 3.1, \uline{on a les isomorphismes suivants}:
@@M29@@


